\section{Discussion}


The mechanism is worth deploying when there is recurrence. The cost saving appears
wherever reuse occurs, but the accuracy penalty is concentrated on UGR'16, where regimes
recur and their label relationship does not. There a reused policy is stale, and the
prospective validation in RAPT-Guard earns its place by taking $0.9137$ to $0.9357$, still
$0.0336$ below the frozen baseline. Reusing on the regime identifier alone is unsafe when
label semantics can drift independently of it, and even a direct validation gate only
repairs most of the damage. The gate is nevertheless cheap: it costs one scoring pass per
regime return, and on the two streams where reuse is safe it matches or beats the frozen
baseline. The practical rule we can support is narrow and honest: reuse when regimes
recur \emph{and} validate the stored policy against the prefix fallback at every return;
if the validation cannot find a stored policy that beats the fallback, do not reuse.


\section{Independent Re-validation}
\label{sec:revalidation}

The primary comparison is re-derived from the raw per-window records by a harness that
re-implements the streaming protocol rather than reading the committed aggregates. It
reproduces Table~\ref{tab:primary} on every predictive metric, and the committed
statistics, gate counts and detector events are all regenerated from those same raw
records. Three findings are worth recording.

\emph{Cost is machine-dependent.} A re-run of the drift pipeline reproduced every F1,
accuracy, precision and recall exactly but shifted adaptation CPU by up to $0.14$\,s. The
cost \emph{ordering} is stable and the gaps large (64--91\%); the absolute seconds are
indicative, not exact.

\emph{Retrieval keys do not track label semantics.} An oracle retraining on the previous
visit of the \emph{same} regime reaches only $0.5926 \pm 0.0077$ macro-F1 on UGR'16,
against $0.9595$ for Full Retraining, but $0.9905 \pm 0.0014$ on Campus. Regime identity
survives on Campus and not on UGR'16, exactly where identity-only reuse collapses.
Recurrence alone does not license reuse: on Campus a same-regime revisit predicts the
previous visit well (mean F1 $0.9876$), yet the class-conditional feature distributions
differ by up to a factor of $10^{8}$.

\emph{Controls.} On Campus a 20-tree cheap refit and a periodic 5-window schedule sit on
the cost--accuracy frontier and neither dominates RAPT-Cheap. The two baselines we add
here separate the trigger from reuse: Frozen+Periodic is the best model on NordicDat
($0.4872$) and Event-Driven+Reuse is the best on UGR'16 ($0.9639$), which shows that on
UGR'16 the event trigger plus reuse recovers what identity-keyed reuse loses, whereas the
RAPT family does not. A label-delay variant, where the target arrives one window late,
leaves Frozen and Event-Driven unchanged but costs RAPT $0.13$ macro-F1 on NordicDat and
Full Retraining $0.06$; the headline numbers assume labels arrive with their window.

\section{Limitations and Threats to Validity}
\label{sec:limitations}
The seed-level tests use thirty seeds, so they are no longer power-limited, but the
mechanism ablation in Table~\ref{tab:ablation} still uses five seeds and the smallest
attainable two-sided $p$-value there is $0.0625$. Window-level tests assume windows from
one stream are independent, which inflates the effective sample size and weakens any claim
meant to generalise. The mechanism relies on regime recurrence, which we measure and do
not assume: a stream with no recurrence never exercises the repository, and a novel regime
falls back to retraining with no advantage over the baseline. The RAPT-Floor threshold and
the 20-tree cheap refit were calibrated for this ensemble; both would need revisiting where
the pre/post-drift gap is larger. NordicDat stays hard (best macro-F1 $0.4872$) and we flag
it as a hard stream rather than a success. Finally, the repository is never pruned, and the
approach treats regime similarity as meaningful, which may fail if the regime signal is
noisy.
\section{Conclusion}
Retraining is the standard response to drift, but when conditions return, part of what was
paid for is already in a previously trained model. RAPT keeps one policy per regime and
retrieves it on recurrence, training anew only when no suitable policy exists. Across four
network streams adaptation CPU was 64--91\% below full retraining, while macro-F1 was lower
on all four. Where regimes recurred with changed label relations, identity-only reuse
dropped below even the frozen baseline; RAPT-Guard, which validates the stored policy
against the prefix fallback before every reuse, recovered most of that loss but on UGR'16
still did not beat the frozen baseline. Policy reuse is worth deploying when regimes recur
and the stored policy is validated at every return, and the natural next step is to test
whether a retrieved regime still means what it did when its policy was stored.
